% ادامه‌ی chapter4.tex

\begin{lessonbox}[جمع‌بندی فصل چهارم]
\begin{itemize}
  \item هندوستان نشان می‌دهد که مدیریت تنوع عظیم در چارچوب دموکراسی فدرال ممکن است — مشروط بر انعطاف‌پذیری نهادی.
  \item بریتانیا اثبات می‌کند که واگذاری نامتقارن می‌تواند بدون تغییر رسمی ماهیت دولت، مطالبات ملّی را مدیریت کند.
  \item سوئیس الگویی آرمانی اما دشوار برای تقلید است: نیازمند فرهنگ سیاسی اجماع‌محور.
  \item بلژیک هشدار می‌دهد که فدرالیسم قومی بدون مکانیسم حل بن‌بست به فلج سیاسی می‌انجامد.
  \item عراق و اتیوپی نشان می‌دهند که طراحی نهادی بدون اعتمادسازی و توسعه‌ی اقتصادی عادلانه بی‌ثمر است.
  \item برای ایران: ترکیبی خلاقانه از عناصر هندی (فدرالیسم نامتقارن)، بریتانیایی (واگذاری تدریجی) و سوئیسی (اجماع‌سازی) قابل تأمل است.
\end{itemize}
\end{lessonbox}